\chapter{Conclusion générale}

GlassInk Monitor transforme des données industrielles dispersées en une chaîne
de supervision cohérente. Le prototype reçoit les messages MQTT, les valide
dans Node-RED, les convertit en Line Protocol, les historise dans InfluxDB et
les présente dans trois tableaux de bord Grafana. Il suit l'état de l'imprimante,
les consommables, les défauts et le résultat des marquages.

Le point le plus délicat n'a pas été l'affichage d'une jauge. Il a été de
définir ce qu'est réellement un arrêt. Sur une ligne qui fonctionne selon les
commandes, une machine immobile n'est pas forcément indisponible. L'ajout d'un
signal de demande rend les indicateurs et les alertes crédibles. La comparaison
entre épuisement et péremption répond au même principe : un pourcentage seul ne
suffit pas pour décider quand intervenir.

Le recours au simulateur a permis d'avancer sans accès au Litmus Edge. Il a
fourni des périodes sans commande, des défauts, des remplissages et des pertes
de communication. Les essais ont validé les doublons MQTT, le journal de rejet,
la reprise après redémarrage, les 23 requêtes de tableaux de bord et les sept
règles d'alerte.

Le projet n'est pas encore raccordé à la production. La prochaine étape est
précise : obtenir un payload Litmus et la liste des tags, confirmer le modèle
de l'imprimante et les unités, puis remplacer le mapping du simulateur. Cette
phase devra aussi valider la sécurité du réseau, le TLS, les ACL, les comptes
Grafana et la sauvegarde.

Ce PFA m'a permis de relier plusieurs domaines de ma formation : réseaux
industriels, messagerie IoT, bases de données temporelles, conteneurisation et
supervision. Il m'a surtout appris qu'un indicateur industriel n'est utile que
si sa définition correspond au fonctionnement réel de la ligne.
